\newcommand{\firmaBox}[2]{%
  \begin{minipage}[t]{4.2cm}
    \centering
    \IfFileExists{firmas/#2.png}%
      {\includegraphics[width=3.8cm,height=1.7cm,keepaspectratio]{firmas/#2.png}}%
      {\IfFileExists{firmas/#2.jpg}%
        {\includegraphics[width=3.8cm,height=1.7cm,keepaspectratio]{firmas/#2.jpg}}%
        {\IfFileExists{firmas/#2.jpeg}%
          {\includegraphics[width=3.8cm,height=1.7cm,keepaspectratio]{firmas/#2.jpeg}}%
          {\fbox{\begin{minipage}[c][1.7cm][c]{3.7cm}\centering\footnotesize\textit{(firma pendiente)}\end{minipage}}}%
        }%
      }\\[2pt]
    \rule{3.8cm}{0.4pt}\\[2pt]
    \small #1
  \end{minipage}%
}

\noindent\textbf{Grupo:} 24

\vspace{0.35cm}

\noindent\textbf{¿Hicieron uso de herramientas de IA en este informe y/o
presentación?} \quad Sí: $\boxtimes$ \qquad No: $\square$

\vspace{0.35cm}

\noindent\textbf{Marquen los usos que correspondan:}

\begin{itemize}
\setlength{\itemsep}{2pt}
\item[$\boxtimes$] Apoyo a la redacción, ortografía o estructura gramatical.
\item[$\square$] Apoyo al análisis de datos e información del problema.
\item[$\square$] Búsqueda o filtro de referencias y literatura.
\item[$\boxtimes$] Apoyo al desarrollo del modelo.
\item[$\boxtimes$] Apoyo al desarrollo de código computacional.
\item[$\boxtimes$] Apoyo al análisis de resultados.
\item[$\boxtimes$] Apoyo para construir figuras o diagramas.
\item[$\boxtimes$] Apoyo para obtener conclusiones.
\item[$\square$] Otro: \rule{8cm}{0.4pt}
\end{itemize}

\noindent\textbf{Herramientas de IA utilizadas y reflexión breve sobre su
aporte, utilidad y verificación de sus resultados:}

\vspace{0.2cm}
Se usaron los modelos Opus y Sonnet 5.5 de Claude principalmente para el desarrollo de códigos computacionales a partir de la modelación. También se usó IA como apoyo para la metodología, análisis de resultados computacionales y escritura del informe, este último sobre todo como etapa final para reducir el número de páginas y corregir la ortografía. Para estos apoyos también se usó GPT-5.6. Estas herramientas fueron útiles para agilizar el desarrollo y mejorar la calidad del trabajo. Sin embargo, sus respuestas y propuestas fueron revisadas por el equipo, verificando la consistencia del código con el modelo formulado y la coherencia de los resultados con las restricciones y objetivos del problema. De esta manera, las herramientas de IA se utilizaron como apoyo, manteniendo la responsabilidad del equipo sobre las decisiones metodológicas y las conclusiones obtenidas.

\vspace{0.5cm}

\noindent\textbf{Nombres y firmas de todos los integrantes:}

\vspace{0.3cm}

\begin{center}
\begin{tabular}{ccc}
\firmaBox{Nicol\'as Bennett}{nicolas_bennett} & \firmaBox{Antonio De Frutos}{antonio_defrutos} & \firmaBox{Claudio Hasb\'un}{claudio_hasbun} \\
\firmaBox{Agust\'in Irarr\'azaval}{agustin_irarrazaval} & \firmaBox{Joaqu\'in Jim\'enez}{joaquin_jimenez} & \firmaBox{Alonso Muelas}{alonso_muelas}
\end{tabular}
\end{center}